\section{Exact Conductor for the Fixed Case \(D=\{3,4,5\}\), \(k=1\)}

Let
\[
B=\{3^a:a\ge 1\}\cup\{4^b:b\ge 1\}\cup\{5^c:c\ge 1\}.
\]
Let \(\operatorname{SubSum}(B)\) denote the set of all finite subset sums of distinct
elements of \(B\).

\begin{theorem}
For the fixed case \(D=\{3,4,5\}\), \(k=1\), the exact conductor is
\[
C(\{3,4,5\},1)=80.
\]
Equivalently,
\[
[80,\infty)\subseteq \operatorname{SubSum}(B)
\]
and
\[
79\notin \operatorname{SubSum}(B).
\]
\end{theorem}

\begin{proof}
Let
\[
F=B\cap[1,2187)
=\{3,4,5,9,16,25,27,64,81,125,243,256,625,729,1024\}.
\]
Then
\[
\sum_{f\in F} f=3236.
\]
The finite replay verifier exhaustively enumerates subset sums of \(F\) and confirms
\[
[80,3156]\subseteq \operatorname{SubSum}(F),
\]
where \(3156=3236-80\). The same endpoint check confirms
\[
79\notin \operatorname{SubSum}(B).
\]
This obstruction is finite because no basis element greater than \(79\) can occur in a
representation of \(79\).

It remains to show the infinite tail. Let \(x\in B\) be a basis element with \(x\ge2187\), and set
\[
T(x)=\sum_{\substack{b\in B\\ b<x}}b.
\]
For each \(q\in\{3,4,5\}\), the sum of \(q\)-powers below \(x\) is at least
\[
\frac{x-q}{q-1}.
\]
Therefore
\[
T(x)\ge
\frac{x-3}{2}+\frac{x-4}{3}+\frac{x-5}{4}
=
\frac{13x-49}{12}.
\]
For \(x\ge1957\),
\[
\frac{13x-49}{12}\ge x+159.
\]
Since every basis element after \(F\) is at least \(2187\), every later basis element
satisfies
\[
T(x)\ge x+159.
\]

Now suppose a finite prefix \(P\subset B\) has total \(T\) and satisfies
\[
[80,T-80]\subseteq \operatorname{SubSum}(P).
\]
If \(x\) is the next basis element and \(x\le T-159\), then the translated interval
\[
x+[80,T-80]=[x+80,x+T-80]
\]
overlaps or touches \([80,T-80]\), because \(x+80\le T-79\). Hence
\[
[80,T+x-80]\subseteq \operatorname{SubSum}(P\cup\{x\}).
\]
The verified prefix \(F\) gives the base case. The inequality above gives the induction
step for every later basis element. Thus
\[
[80,\infty)\subseteq \operatorname{SubSum}(B).
\]
Together with \(79\notin\operatorname{SubSum}(B)\), this proves
\[
C(\{3,4,5\},1)=80.
\]
\end{proof}
